\section{求导规则从哪里来}
\label{sec:02-Calculus/03-Differentiation/02}

按定义求 \(f(x)=x^3\sin x+e^x\cos x\) 的导数，第一步就要写下：
\[
\frac{(x+h)^3\sin(x+h)+e^{x+h}\cos(x+h)-x^3\sin x-e^x\cos x}{h}.
\]
若强行展开和角公式与二项式，极其繁琐且容易出错。

问题的出路在于换个视角。上一节我们将可微性表达为线性近似：
\[
f(x+h)=f(x)+f'(x)h+o(h),
\]
即函数在一点附近可以由一个线性模型逼近，导数就是这个线性模型中 \(h\) 的一次项系数。

因此，求导规则的本质是：\textbf{已知 \(f\) 和 \(g\) 各自的线性模型，求它们经过加、乘、除运算后，新函数的线性模型是什么}。只需三步——代入模型、展开式子、舍去高阶无穷小 \(o(h)\)，剩下 \(h\) 的一次项系数即为所求。

\subsection{加法与乘法法则}

线性组合的推导最为直接。将 $af+bg$ 的线性模型直接相加， $h$ 的系数即为 $af'+bg'$，因此：
\[
(af+bg)' = af'+bg'.
\]

对于乘法，直觉容易误以为 \((fg)'=f'g'\)。但通过模型相乘便能清楚看清原因：
\[
\begin{aligned}
(fg)(x+h) &= \bigl(f+f'h+o(h)\bigr)\bigl(g+g'h+o(h)\bigr) \\
&= fg+(f'g+fg')h+f'g'h^2+o(h).
\end{aligned}
\]
右端第三项 \(f'g'h^2\) 除以 \(h\) 后仍趋于 \(0\)，属于高阶无穷小 \(o(h)\)。保留一阶项系数，即得乘积法则：
\[
(fg)'=f'g+fg'.
\]

几何上的面积变化清晰地展示了这一过程（见下图）：

\begin{center}
\begin{tikzpicture}[x=1cm,y=1cm,font=\footnotesize]
  \fill[blue!12] (0,0) rectangle (3,2);
  \fill[orange!30] (3,0) rectangle (3.9,2);
  \fill[green!25] (0,2) rectangle (3,2.6);
  \fill[red!35] (3,2) rectangle (3.9,2.6);
  \draw[black!60] (0,0) rectangle (3.9,2.6);
  \draw[black!60] (3,0) -- (3,2.6);
  \draw[black!60] (0,2) -- (3.9,2);
  \node at (1.5,1) {\(fg\)};
  \node at (3.45,1) {\(g \cdot (f'h)\)};
  \node at (1.5,2.3) {\(f \cdot (g'h)\)};
  \draw[<->] (0,-0.35) -- (3,-0.35) node[midway,below] {\(f\)};
  \draw[<->] (3,-0.35) -- (3.9,-0.35) node[midway,below] {\(f'h\)};
  \draw[<->] (-0.4,0) -- (-0.4,2) node[midway,left] {\(g\)};
  \draw[<->] (-0.4,2) -- (-0.4,2.6) node[midway,left] {\(g'h\)};
  \draw[black!60] (3.9,2.3) -- (4.4,2.3);
  \node[right] at (4.4,2.3) {\(f'g'h^2\)：高阶微元，省略};
\end{tikzpicture}

\emph{矩形面积的增量由三部分组成：边长单侧增加产生的两个长条（一阶微元），以及右上角由两侧同时增加产生的小方块（二阶微元 \(h^2\)）。在一阶尺度下，仅需保留两个长条的贡献。}
\end{center}

有了加法与乘法法则，开头提到的复杂函数便可直接拆解求导：
\[
\bigl(x^3\sin x+e^x\cos x\bigr)'=(3x^2\sin x+x^3\cos x)+(e^x\cos x-e^x\sin x).
\]

\subsection{商法则的简化推导}

商法则无需另起炉灶。先考虑倒数 $1/g$ 的线性近似，利用代数通分：
\[
\frac{1}{g(x+h)} - \frac{1}{g(x)} = \frac{1}{g + g'h + o(h)} - \frac{1}{g} = \frac{-g'h + o(h)}{g(g + g'h + o(h))} = -\frac{g'}{g^2}h + o(h).
\]
由此可知：
\[
\left(\frac{1}{g}\right)' = -\frac{g'}{g^2}.
\]
再结合乘积法则处理 $f \cdot (1/g)$，即可顺理成章地导出商法则：
\[
\left(\frac{f}{g}\right)' = \frac{f'g - fg'}{g^2}.
\]
分母非零不仅是公式成立的免责声明，更是 $1/g$ 在该点处能够建立线性模型的前提。

\subsection{基本函数的导数源于极限定义}

运算法则解决了复杂函数的组装问题，但最基础的构建模块仍需由极限定义提供。

{\def\LTcaptype{none}
\begin{longtable}{@{}lll@{}}
\toprule
函数 & 导数 & 核心推导依据 \\
\midrule
\(x^n\) & \(nx^{n-1}\) & 二项式展开的一阶项系数 \\
\(\sin x\) & \(\cos x\) & 重要极限 \(\lim_{h\to0}\frac{\sin h}{h}=1\)（要求必须使用弧度制） \\
\(\cos x\) & \(-\sin x\) & 同上 \\
\(e^x\) & \(e^x\) & 自然底数 \(e\) 的定义使得 \(\lim_{h\to0}\frac{e^h-1}{h}=1\) \\
\(\ln x\) & \(1/x\quad (x>0)\) & 指数函数 \(e^x\) 的反函数性质 \\
\(a^x\) & \(a^x\ln a\) & 改写为恒等式 \(e^{x\ln a}\) 展开 \\
\bottomrule
\end{longtable}
}

明确每个基本公式的支撑点至关重要：若单位换为角度制，正弦函数的导数会多出 \(\pi/180\) 的系数；若改变底数，指数函数的导数则会多出 \(\ln a\) 的缩放因子。

\subsection{非标准形式的改写技巧}

当遇到的表达式无法直接套用基本规则时（例如 \(y=x^x\)），通常不需要引入新法则，而是通过代数变换将其变形为已知规则管辖的形状。

例如对 \(y=x^x\quad (x>0)\) 两边取对数：
\[
\ln y = x\ln x.
\]
对两侧同时求导，左侧按隐函数/复合函数求导得 \(\frac{y'}{y}\)，右侧按乘积法则得 \(\ln x+1\)，整理即得：
\[
y' = x^x(\ln x+1).
\]
无论是将根式化为分数指数幂、将连乘转化为对数求和，本质都是利用等价变换将式子纳入既有的规则体系之中。

\subsection{最关键的构造：复合与嵌套}

至此，代数加减、乘法与除法三种组装方式均已处理完毕。然而在实际应用中，还有一种至关重要的函数构造方式——\textbf{函数的嵌套}（如 \(\sin(x^2)\)）。

嵌套结构中，自变量的变化需要经过中间层才能传递至外层。处理这种链式传递的规则（复合函数求导法则），其重要性远超基本的代数运算规则：深层神经网络本质上就是由多层线性与非线性函数嵌套而成的庞大系统，而训练此类模型所依赖的梯度反向传播算法，核心正是这一条复合求导规则。

下一节我们将专门讨论它。